\begin{table*}[h!]
\small
\caption{{\ben} Book links \& Statistics}
\label{table:benchmark}
\resizebox{\textwidth}{!}{%
% \begin{center}
\begin{tabular}{llllllll}
\textbf{Domain} & \textbf{Book} & \textbf{\#Pages} & \textbf{\#Leaves} &\textbf{Test}&\textbf{Train}&\textbf{Dev}\\
\toprule
\multirow{3}{*}{Education} & \href{https://open.umn.edu/opentextbooks/textbooks/1254}{Open Music Theory} & 1297 & 429 & 7430 & 4493 & 1671  \\
& \href{https://open.umn.edu/opentextbooks/textbooks/1621}{The Whole Child: Development in the Early Years} & 182 & 129 & 1746 & 1056& 388 \\
&  \href{https://open.umn.edu/opentextbooks/textbooks/221}{Teaching in a Digital Age} & 779 & 118 & 6150&3797&1489\\
\midrule
\multirow{3}{*}{Finance} & \href{https://open.umn.edu/opentextbooks/textbooks/principles-of-financial-accounting}{Principles of Financial Accounting} & 318 & 98 & 4119&2395&981 \\
& \href{https://open.umn.edu/opentextbooks/textbooks/financial-accounting}{Accounting in the Finance World} & 572 & 80 & 4481&2747&1081 \\
& \href{https://open.umn.edu/opentextbooks/textbooks/financial-and-managerial-accounting}{Financial and Managerial Accounting} & 1077 & 107 & 7815&4692&1901\\
\midrule
\multirow{3}{*}{Law} & \href{https://open.umn.edu/opentextbooks/textbooks/1219}{Construction Contracting} & 403 & 103 & 3653&2249&866\\
& \href{https://open.umn.edu/opentextbooks/textbooks/criminal-procedure-undergraduate-edition-author}{Criminal Procedure} & 897 & 106 & 5610& 3401& 1349 \\
& \href{https://open.umn.edu/opentextbooks/textbooks/tort-law-cases-and-commentaries}{Tort Law: Cases and Commentaries} & 948 & 396 & 18209&11174&4363 \\
\midrule
\multirow{3}{*}{Medicine} & \href{https://open.umn.edu/opentextbooks/textbooks/1321}{Nursing Assistant} & 659 & 118 & 4586&2790&1093 \\
& \href{https://open.umn.edu/opentextbooks/textbooks/1013}{Nursing Fundamentals} & 1327 & 121 & 9065&5493&2200 \\
& \href{https://open.umn.edu/opentextbooks/textbooks/1427}{Nursing Management and Professional Concepts} & 599 & 74 & 3695&2277&885 \\
\midrule
\multirow{3}{*}{Natural Sciences} & \href{https://open.umn.edu/opentextbooks/textbooks/1279}{Introduction to Genetics} & 513 & 70 & 2254&1345&527 \\
& \href{https://open.umn.edu/opentextbooks/textbooks/1600}{Principles of Mechanics} & 179 & 125 & 2341&1418&543 \\
& \href{https://open.umn.edu/opentextbooks/textbooks/organic-chemistry}{Organic Chemistry} & 1249 & 321 & 11124&6601&2635 \\
\midrule
\multirow{3}{*}{Social Sciences} & \href{https://open.umn.edu/opentextbooks/textbooks/1446}{Foundations of Aural Skills} & 674 & 118 & 2451&1421&559 \\
& \href{https://open.umn.edu/opentextbooks/textbooks/1435}{Introduction to Comparative Government and Politics} & 421 & 189 & 3893&2288&898 \\
& \href{https://open.umn.edu/opentextbooks/textbooks/a-practicum-in-behavioral-economics}{A Practicum in Behavioral Economics} & 381 & 178 & 4627&2821&1073 \\
\end{tabular}
% \end{center}
}
\end{table*}

\subsection{Breaking the Mold: Why a New Benchmark Is Needed}
There are multiple benchmarks proposed around long context applications such as ContractNLI~\citep{koreeda-manning-2021-contractnli-dataset} focused on NLI or \textbf{Scrolls~\citep{shaham-etal-2022-scrolls}} with seven challenging tasks. However, there doesn't exist any benchmark for long context retrieval with any form of structured view, which we hypothesize as the key for precise retrieval. Benchmarks such as GovReport~\citep{huang2021govreport}, SummScreen~\citep{chen-etal-2022-summscreen}, and QM-Sum ~\cite{zhong2021qmsum} focus on generating summaries of a length document, whereas, Qasper ~\citep{dasigi2021qasper}, QuALITY~\citep{pang-etal-2022-quality} and NarrativeQA~\citep{kocisky2018narrativeqa} focus on the generation aspect with \textbf{only answers} for a user query and \textbf{do not contain the gold passage for retrieval}. The closest benchmark is \textbf{LocoV1~\citep{SaadFalcon2024BenchmarkingAB}} which does include gold passages for each query. However, due to \textbf{lack of {\toc} or a global structure for the input document}, we cannot use this benchmark for our task.\\ To the best of our knowledge, {\ben} is the first benchmark which provides a clean structured view with Table of Contents for evaluating long context retrieval with structures . 

\subsection{{\ben}: A diverse and novel benchmark for {\toc} Retrieval}

Our main motivation for a new benchmark is to evaluate how well a global semantic structure such as a {\toc} can help in better retrieval for long context applications which is typical for technical reports or voluminous books. Thus to create a clean and effective benchmark we turn to \href{https://open.umn.edu/opentextbooks/}{opentextbooks}, which is perhaps the largest collection of such textbooks. To make the benchmark a comprehensive one across domains, we picked six diverse domains namely -- Education, Finance, Law, Medicine, Natural Sciences and Social Sciences and selected three books from each domain (Table~\ref{table:benchmark}). We parsed the PDF for each book and extracted Table of Contents using pymupdf\footnote{\url{https://pymupdf.readthedocs.io/en/latest/}}. A sample of the final cleaned content for one of the books can be viewed \href{https://anonymous.4open.science/r/s_331/data/education/book1/v1/content.jsonl}{here}. For each paragraph we ask a powerful LLM Mixtral 8x7b model~\citep{jiang2024mixtral} to generate multiple questions covering all important topics in the paragraph~\cite{Zhang2024SyntheticKI}. Following the same technique as was proposed in DSI~\cite{Tay2022TransformerMA} to train the parametric index, we use a portion of these questions for training DSI~\cite{Tay2022TransformerMA} and {\sys}, a small portion as development set which is to help us pick the best checkpoint. Majority of the generated questions were picked as test questions -- to thoroughly test coverage of content by an IR system. We believe our benchmark {\ben} will further research in building IR systems which can leverage structure in a corpus and retrieve the most relevant chunks for a user query.



% We take special care in creating the test set so that it doesn't suffer from LLM bias. We adopt two complementary strategies to make the test set more robust: 
% (1) Through LLM based generation: We use the LLM to generate test queries over the paragraphs, but with changes in generation strategies through temperature scaling as well with different prompts to minimize overlap with train data.  We do ensure to have enough examples to thoroughly test the full coverage of the long content by an IR system. 
% (2) Through human generation: We also adopt a complementary approach of human-in-the-loop test set generation which closely mimics how a teacher (i.e. the human) would test a student (i.e. {\sys}), whether she understands the book (i.e. the search corpus) well enough by testing her ability to recall the relevant content (i.e. \toc~ entry) given a query. We use a working prototype of {\sys} and ask $5$ human annotators working on $5$ different books across $5$ domains to open up any random paragraph from a particular book and frame a query by looking at the paragraph content. The annotators try these human generated queries in {\sys} and check whether {\sys} was able to retrieve the right \toc~ entry. 

% We add both of these generated queries to the benchmark. Because human generated queries have no scope of LLM bias, we use the results from this human test set and try to correlate it with the results and observations obtained from test set with LLM based generations. We ensure they corroborate each other to nullify scope of any LLM bias in designing {\ben}. We believe our benchmark {\ben} will further research in building IR systems which can leverage structure in a corpus and retrieve the most relevant chunks for a user query.